%
\againframe{plan}
%
\begin{frame}{Lecture plan}
	\begin{learningobjectives}
		\begin{itemize}
			\item Condensed matter at macroscopic scale
			\item Classification of condensed matter based on conductivity
			\item Nature of metals
			\item Classical electron theory of metals -- Assumptions
			\item Ohm's law for metals -- Conductivity
			\item Application of metallic conductivity
		\end{itemize}
	\end{learningobjectives}
\end{frame}
\begin{frame}{Discovery of material}
	\begin{figure}
		\centering
		\includegraphics[height=0.6\textheight]{ages_human_civilization_spacetimeline}
		\caption{Civilization spacetimeline: Copper \(\rightarrow\) Bronze \(\rightarrow\) Iron \(\rightarrow\) Steel \(\rightarrow\) Silicon}
	\end{figure}
	%
	\begin{keyinsight}
		Material defines the age.
	\end{keyinsight}
\end{frame}
%
\begin{frame}{Condensed matter}
	\begin{itemize}
		\item Material in the liquid or solid form is called \textbf{condensed matter}.
		\item Condensed matter is further sub-classified based on  electrical, optical, magnetic, thermal, mechanical properties at the \textbf{macroscopic} scale.  In the case of electrical property, we apply electric field and classify the materials based on their conductivity.
		\item The macroscopic behaviour is related to the \textbf{microscopic} behaviour of electrons under applied ``forces''.
	\end{itemize}
	\begin{learningobjectives}
		To relate the macroscopic properties with the microscopic behaviour of electrons in condensed matter.
	\end{learningobjectives}
\end{frame}
%
\begin{frame}{Classification of Condensed Matter by Conductivity}
	%
	\begin{figure}
		\centering
		\includegraphics[width=\textwidth]{condensed_matter_electrical_classification}
		\caption{(a) Gold, a metal, (b) Silicon, a semiconductor, (c) Diamond, an insulator}
	\end{figure}
	%
	\begin{minipage}{0.85\textwidth}
		\begin{itemize}
			\item Conductivity is the measure of how easily electrons move under applied electric field. Its unit is \si{\per\ohm\per\metre} or \si{\siemens\per\metre} (S for Siemens).
			%
			\item Materials can be classified based on conductivity as:
			\begin{enumerate}
				\item \textbf{Metals}: High ($\sigma \sim 10^7$ S/m).
				\item \textbf{Semiconductors}: Intermediate ($\sigma \sim 10^{-4}$ S/m).
				\item \textbf{Insulators}: Negligible ($\sigma \sim 10^{-10}$ S/m).
			\end{enumerate}
		\end{itemize}
	\end{minipage}
	%
	\begin{minipage}{0.13\textwidth}
		\begin{figure}
			\centering
			\includegraphics[width=\textwidth]{Werner_Siemens_1816-1892}
			\caption{\centering \footnotesize W. Siemens [1816-1892]}
		\end{figure}
	\end{minipage}
	
\end{frame}
%
\begin{frame}{Elemental phases}
	\begin{itemize}
		\item The electrical state of condensed matter is also called a \textbf{phase} -- similar to solid phase, liquid phase, etc.
	\end{itemize}
	\begin{figure}
		\centering
		\includegraphics[height=0.5\textheight]{condensed_matter_periodic_table}
	\end{figure}
	%
	\begin{keyinsight}
		Most elemental phases are metals.
	\end{keyinsight}
\end{frame}
%
%
\begin{frame}{Conductivity of phases}
	%
	\begin{itemize}
		\item Conductivity \(\sigma\) is inversely related to resistivity \(\rho\) by \(\rho = \tfrac{1}{\sigma}\)
	\end{itemize}	
	%
	\begin{minipage}{0.75\textwidth}
		\begin{figure}
			\centering
			\includegraphics[width=\textwidth]{metals_resistivity_chart}
			%	\caption{Resistivity chart for metals, semiconductor and insulators}
		\end{figure}
	\end{minipage}
	%
	\hfill
	%
	\begin{minipage}{0.23\textwidth}
		{
			\footnotesize	
			\begin{table}
				\centering
				\begin{tabular}{l l }
					\toprule
					&		\(\rho\)	\\
					&		\si{\ohm\metre} \\
					\midrule
					\ce{Cu}	 &	\num{e-7}	\\
					\ce{Si} 	&	\num{e4}	\\
					\ce{SiO2}&	\num{e10}	\\
					\bottomrule
				\end{tabular}
			\end{table}
		}
	\end{minipage}
	%
	\begin{keyinsight}
		Conductivity spans ``orders of magnitude'' across phases.
	\end{keyinsight}
	%
\end{frame}
%
\begin{frame}{Early experiments:Avogadro number and Electron charge}
		\begin{figure}
		\includegraphics[width=2.2\textwidth]{condensed_macroscopic_to_microscopic}
	\end{figure}
	%
\end{frame}
%
\begin{frame}{Macroscopic \(\rightarrow\) microscopic}
	%
	\begin{figure}
		\includegraphics[width=\textwidth]{condensed_macroscopic_to_microscopic}
	\end{figure}
	%
	\begin{columns}
		\begin{column}[t]{0.7\textwidth}
			%
			\begin{guesstimate}{Avogadro number \(N_A\)}
				To electroplate \SI{63.5}{\gram} of copper, it takes \SI{2}{\faraday} of charge. [Hint: 1F (F for Faraday)= \SI{96,485}{\coulomb}, charge of electron \(e\) = \SI{1.602e-19}{\coulomb}]
			\end{guesstimate}
			%
			\begin{guesstimate}{Radius of atom}
				The density of copper is \SI{8.96}{\gram\per\centi\metre\cubed}.
			\end{guesstimate}
		\end{column}
		%
		\begin{column}[t]{0.28\textwidth}
			%
			\begin{figure}
				% Base image
				\includegraphics[width=0.48\textwidth]{Michael_Faraday_1791_1867}
				\begin{overpic}[width=0.48\textwidth]{Robert_Millikan_1868_1953}
					% Overlay image on top
					\put(0,0){\includegraphics[width=0.2\textwidth]{Alfred_Nobel_1833_1896}}
				\end{overpic}
				\caption{Faraday, Millikan}
			\end{figure}
		\end{column}
	\end{columns}
\end{frame}
%
